\subsection{The Philosophy of Syntactic Whitespace}

In many programming languages, whitespace is mostly decorative. A C program may be written with careful indentation, ugly indentation, or almost no indentation, and the compiler still recognizes the structure from explicit tokens such as \texttt{\{} and \texttt{\}}. Python makes a different choice. In Python, indentation is not merely a visual aid. It is part of the syntax.

This design often feels strange to programmers trained in C, because C separates visual layout from formal block structure. Python deliberately refuses that separation. The visual structure and the executable structure are forced to agree.

\begin{quote}
In Python, indentation is not a comment made by the programmer to the reader. It is structural information consumed by the interpreter.
\end{quote}

This choice connects directly to the previous control-flow discussion. An \texttt{if} statement, a \texttt{while} loop, a \texttt{for} loop, a \texttt{try} block, an \texttt{except} block, and a function body all need a way to mark which statements belong inside them. Python uses indentation for that boundary.

\subsubsection{Python’s Core Design Choice: Whitespace Semantics for Structural Block Definitions}

A block is a group of statements controlled by another statement. In Python, the block starts after a header line ending in a colon. The indented statements below that header belong to the block.

\begin{verbatim}
mood = "confident"

if mood == "confident":
    print("Nobody expects the Spanish Inquisition!")
    print("This line belongs to the if block.")

print("This line is outside the if block.")
\end{verbatim}

Possible output:

\begin{verbatim}
Nobody expects the Spanish Inquisition!
This line belongs to the if block.
This line is outside the if block.
\end{verbatim}

The indentation is not optional. The two indented \texttt{print()} statements are controlled by the \texttt{if}. The final \texttt{print()} statement is not controlled by the \texttt{if}, because it returns to the previous indentation level.

The control-flow structure is therefore visible directly in the left margin:

\begin{verbatim}
if condition:
    inside block
    inside block
outside block
\end{verbatim}

The colon announces that a nested block is expected. The increased indentation defines the body of that block. Returning to the earlier indentation level ends the block.

This means that Python code cannot lie about its structure in the same way badly formatted C code can. In C, this is legal but misleading:

\begin{verbatim}
/* C example, not Python */
if (ready)
    printf("ready\n");
    printf("still printed even when ready is false\n");
\end{verbatim}

The second \texttt{printf()} visually looks connected to the \texttt{if}, but it is not controlled by the \texttt{if}. In Python, that ambiguity is removed because indentation itself defines the block.

\begin{verbatim}
ready = False

if ready:
    print("ready")

print("still printed even when ready is false")
\end{verbatim}

Possible output:

\begin{verbatim}
still printed even when ready is false
\end{verbatim}

If the second \texttt{print()} should belong to the \texttt{if}, it must be indented:

\begin{verbatim}
ready = False

if ready:
    print("ready")
    print("also controlled by the if")
\end{verbatim}

No output appears, because both \texttt{print()} statements are inside the false branch.

The rule is direct:

\begin{quote}
Same indentation level means same structural layer. Greater indentation means a nested block. Returning to an earlier indentation level closes one or more nested blocks.
\end{quote}

\subsubsection{Structural Isolation: Eliminating Explicit Block Tokens (Curly Braces \texttt{\{\}}, \texttt{begin}/\texttt{end})}

Python does not use curly braces to define ordinary code blocks. It also does not use paired words such as \texttt{begin} and \texttt{end}. Instead, the block boundary is carried by indentation.

A C-style structure uses explicit delimiters:

\begin{verbatim}
/* C example, not Python */
if (score == 42) {
    printf("the answer\n");
    printf("do not panic\n");
}
printf("done\n");
\end{verbatim}

A Python structure uses colon plus indentation:

\begin{verbatim}
score = 42

if score == 42:
    print("the answer")
    print("do not panic")

print("done")
\end{verbatim}

Possible output:

\begin{verbatim}
the answer
do not panic
done
\end{verbatim}

The missing curly braces are not an omission. They are a design decision. Python removes one class of structural tokens and transfers their job to the layout of the code.

This has several consequences.

First, block structure becomes visually unavoidable. If two statements are aligned, Python treats them as belonging to the same block level.

\begin{verbatim}
name = "Kramer"

if name == "Kramer":
    print("Giddy up!")
    print("This is still inside the if.")

print("This is outside the if.")
\end{verbatim}

Second, indentation mistakes become syntax mistakes, not merely style mistakes. If Python expects an indented block and does not receive one, the program is invalid.

\begin{verbatim}
name = "Bart"

if name == "Bart":
print("Eat my shorts.")
\end{verbatim}

This raises an indentation error before the program can run normally:

\begin{verbatim}
IndentationError: expected an indented block after 'if' statement
\end{verbatim}

Third, Python makes the physical shape of the code part of the language contract. This does not mean that any whitespace matters. Spaces between words and operators are usually flexible. The special whitespace is the indentation at the beginning of logical lines.

For example, these assignments are equivalent:

\begin{verbatim}
x=42
x = 42
x     =     42
\end{verbatim}

But these block structures are not equivalent:

\begin{verbatim}
if x == 42:
    print("inside")

print("outside")
\end{verbatim}

and:

\begin{verbatim}
if x == 42:
    print("inside")
    print("also inside")
\end{verbatim}

The left margin changes the control-flow graph.

\subsubsection{Lexer Parsing Rules: How the Tokenizer Tracks Indentation via Stacked \texttt{INDENT} / \texttt{DEDENT} States and Generates \texttt{IndentationError}}

Python does not wait until execution to discover block structure. The source text is first processed by the tokenizer. The tokenizer reads logical lines and emits tokens. Among these tokens are special structural markers named \texttt{INDENT} and \texttt{DEDENT}.

A useful mental model is an indentation stack.

At the start of a file, Python is at indentation level zero. When a new block begins and the next logical line is more indented than the previous block level, Python pushes the new indentation level and emits an \texttt{INDENT} token. When a later line returns to an earlier indentation level, Python pops indentation levels and emits one or more \texttt{DEDENT} tokens.

The following script shows these tokens using the standard \texttt{tokenize} module:

\begin{verbatim}
import io
import tokenize

src = """\
score = 42

if score == 42:
    print("the answer")
    print("do not panic")

print("done")
"""

tokens = tokenize.generate_tokens(io.StringIO(src).readline)

for tok in tokens:
    tok_name = tokenize.tok_name[tok.type]
    if tok_name in ("INDENT", "DEDENT", "NAME", "OP", "STRING", "NUMBER"):
        print(f"{tok_name:<8} {tok.string!r}")
\end{verbatim}

Possible output:

\begin{verbatim}
NAME     'score'
OP       '='
NUMBER   '42'
NAME     'if'
NAME     'score'
OP       '=='
NUMBER   '42'
OP       ':'
INDENT   '    '
NAME     'print'
OP       '('
STRING   '"the answer"'
OP       ')'
NAME     'print'
OP       '('
STRING   '"do not panic"'
OP       ')'
DEDENT   ''
NAME     'print'
OP       '('
STRING   '"done"'
OP       ')'
\end{verbatim}

The important lines are:

\begin{verbatim}
INDENT   '    '
DEDENT   ''
\end{verbatim}

They show that indentation is turned into formal structure before the program executes. The interpreter is not guessing from appearance. The tokenizer has already converted indentation changes into tokens.

The token objects produced by \texttt{tokenize} are runtime objects too:

\begin{verbatim}
import io
import tokenize

src = "x = 42\n"
tokens = tokenize.generate_tokens(io.StringIO(src).readline)

first_tok = next(tokens)

print(f"type(first_tok): {type(first_tok)}")
print(f"first_tok: {first_tok}")
print(f"first_tok.type: {first_tok.type}")
print(f"first_tok.string: {first_tok.string!r}")
print(f"has start: {'start' in dir(first_tok)}")
print(f"has end: {'end' in dir(first_tok)}")
print(f"has line: {'line' in dir(first_tok)}")
\end{verbatim}

Possible output:

\begin{verbatim}
type(first_tok): <class 'tokenize.TokenInfo'>
first_tok: TokenInfo(type=1 (NAME), string='x', start=(1, 0), end=(1, 1), line='x = 42\n')
first_tok.type: 1
first_tok.string: 'x'
has start: True
has end: True
has line: True
\end{verbatim}

The \texttt{type} field identifies the token category. The \texttt{string} field stores the exact source text for that token. The \texttt{start} and \texttt{end} fields store source positions. The \texttt{line} field stores the physical line from which the token came.

If Python expects indentation and the indentation does not form a valid block, it raises \texttt{IndentationError}.

\begin{verbatim}
bad_src = """\
if True:
print("Ni!")
"""

try:
    compile(bad_src, "<bad_indent>", "exec")
except IndentationError as err:
    print(f"type(err): {type(err)}")
    print(f"message: {err}")
    print(f"args: {err.args}")
    print(f"has lineno: {'lineno' in dir(err)}")
    print(f"line number: {err.lineno}")
\end{verbatim}

Possible output:

\begin{verbatim}
type(err): <class 'IndentationError'>
message: expected an indented block after 'if' statement on line 1 (<bad_indent>, line 2)
args: ("expected an indented block after 'if' statement on line 1", ('<bad_indent>', 2, 1, 'print("Ni!")\n', 2, -1))
has lineno: True
line number: 2
\end{verbatim}

This example uses \texttt{compile()} so the error can be caught and inspected instead of stopping the whole demonstration script. The bad source is never executed. It fails during compilation, while Python is building the executable structure.

\subsubsection{Mixed Tabs and Spaces: How \texttt{TabError} Emerges from Ambiguous Indentation}

Python allows indentation with spaces or tabs in principle, but mixing them in the same block structure is dangerous. A tab is not visually stable across editors. One editor may display it as four columns, another as eight columns, and another may use a configurable width.

Python must decide block structure mechanically, not visually. When tabs and spaces create ambiguous indentation, Python raises \texttt{TabError}. \texttt{TabError} is a specialized form of \texttt{IndentationError}.

The following script builds problematic source text as a string. The \texttt{\textbackslash t} sequence represents a tab character.

\begin{verbatim}
bad_src = "if True:\n        print('spaces')\n\tprint('tab')\n"

try:
    compile(bad_src, "<mixed_tabs_spaces>", "exec")
except TabError as err:
    print(f"type(err): {type(err)}")
    print(f"is IndentationError: {isinstance(err, IndentationError)}")
    print(f"message: {err}")
    print(f"args: {err.args}")
\end{verbatim}

Possible output:

\begin{verbatim}
type(err): <class 'TabError'>
is IndentationError: True
message: inconsistent use of tabs and spaces in indentation (<mixed_tabs_spaces>, line 3)
args: ('inconsistent use of tabs and spaces in indentation', ('<mixed_tabs_spaces>', 3, 1, "\tprint('tab')\n", 3, 0))
\end{verbatim}

The line with spaces and the line with a tab may look similarly indented in some editors, but Python rejects the structure because the indentation is not reliably defined.

The inheritance relationship is visible directly:

\begin{verbatim}
print(TabError)
print(TabError.__mro__)
\end{verbatim}

Possible output:

\begin{verbatim}
<class 'TabError'>
(<class 'TabError'>, <class 'IndentationError'>, <class 'SyntaxError'>, <class 'Exception'>, <class 'BaseException'>, <class 'object'>)
\end{verbatim}

This means a handler for \texttt{IndentationError} can also catch \texttt{TabError}, because \texttt{TabError} is lower in the same exception class tree.

\begin{verbatim}
bad_src = "if True:\n        print('spaces')\n\tprint('tab')\n"

try:
    compile(bad_src, "<mixed_tabs_spaces>", "exec")
except IndentationError as err:
    print(f"caught as indentation problem: {type(err)}")
\end{verbatim}

Possible output:

\begin{verbatim}
caught as indentation problem: <class 'TabError'>
\end{verbatim}

The practical rule is simple: use spaces for indentation, and use the same number of spaces for the same block depth. Four spaces per indentation level is the common Python convention.

\begin{quote}
Tabs and spaces are not morally different characters. The problem is that mixed indentation can make the visual block and the tokenizer's block disagree.
\end{quote}

For this course, the deeper point is not merely style. It is structural determinism. Python must transform source text into a definite execution graph. Indentation is one of the inputs to that transformation. If the indentation cannot be converted into a clear nesting structure, Python refuses to compile the program.